\section{Multimodal features in fashion recommendation}

\subsection{Visual Features}

In fashion recommendation, visual attributes (such as color, pattern, style, shape, and fabric texture) are critical determinants of user preferences. Because clothing decisions are heavily guided by aesthetics, purely collaborative signals or textual descriptions often fail to capture these stylistic nuances. Consequently, incorporating visual features is essential to enrich item representations \parencite{10.1007/978-981-95-3355-8_6}. In this study, visual features are processed and stored as \texttt{image\_embeddings.npy} using two alternative encoders:
\begin{itemize}
    \item \textit{MobileNetV2}: A lightweight convolutional neural network pre-trained on ImageNet. It processes the item image \(I_i\) to extract a spatial feature representation:
    \begin{equation}
        \mathbf{x}_i^{\text{vis}} = \text{MobileNetV2}(I_i) \in \mathbb{R}^{d_v},
    \end{equation}
    where \(d_v\) is the visual feature dimension. MobileNetV2 is highly computationally efficient, making it suitable for low-latency inference and resource-constrained deployments, although it lacks deep semantic alignment with textual concepts.
    \item \textit{CLIP (Contrastive Language-Image Pretraining)}: A state-of-the-art vision-language model trained on web-scale image-text pairs. The visual encoder of CLIP maps the image \(I_i\) into a joint vision-language latent space:
    \begin{equation}
        \mathbf{x}_i^{\text{vis}} = \text{CLIP}_{\text{vis}}(I_i) \in \mathbb{R}^{d_v}.
    \end{equation}
    While computationally heavier than MobileNetV2, CLIP features capture high-level semantic abstractions, representing aesthetic style and textual compatibility in a unified space.
\end{itemize}
The trade-off between these encoders is clear: CLIP provides richer semantic representations at the cost of higher feature extraction overhead, whereas MobileNetV2 offers a lightweight and fast alternative for standard convolutional feature extraction.


\subsection{Textual Features}

Textual metadata, such as product titles, descriptions, categories, and tags, provides explicit, descriptive information about fashion items. In this study, text features are extracted and stored as \texttt{text\_embeddings.npy} using two distinct modeling paradigms:
\begin{itemize}
    \item \textit{TF-IDF (Term Frequency-Inverse Document Frequency)}: A statistical approach that constructs sparse, high-dimensional lexical vectors based on term occurrence. For item description text \(T_i\), the TF-IDF representation \(\mathbf{x}_i^{\text{txt}} \in \mathbb{R}^{V}\) (where \(V\) is the vocabulary size) is effective for capturing exact keyword matches and constructing computationally efficient item-item text similarity matrices.
    \item \textit{BERT (Bidirectional Encoder Representations from Transformers)}: A pre-trained transformer model that captures contextual semantic meanings. The BERT encoder processes the text \(T_i\) to output a dense semantic embedding:
    \begin{equation}
        \mathbf{x}_i^{\text{txt}} = \text{BERT}(T_i) \in \mathbb{R}^{d_t},
    \end{equation}
    where \(d_t\) is the text embedding dimension. BERT effectively handles synonyms and semantic similarity, although it is computationally intensive and can suffer from domain shift when applied to specialized e-commerce fashion vocabularies.
\end{itemize}
In this study, TF-IDF and BERT are compared as alternative textual encoders to analyze the impact of lexical matches versus deep semantic representations on recommendation quality.


\subsection{Multimodal Fusion}

To leverage both visual aesthetics and textual attributes, multimodal recommender systems employ fusion strategies to integrate features from different modalities into a unified item representation. In the proposed pipeline, the choice of fusion is parameterized by the \texttt{sim\_type} variable, evaluating the following configurations:
\begin{itemize}
    \item \textit{Visual-only (\texttt{img\_only})}: The item representation is constructed solely from the visual embedding: \(\mathbf{x}_i = \mathbf{x}_i^{\text{vis}}\).
    \item \textit{Textual-only (\texttt{tfidf} or \texttt{bert})}: The item representation is constructed solely from the textual embedding: \(\mathbf{x}_i = \mathbf{x}_i^{\text{txt}}\).
    \item \textit{Late Fusion (\texttt{multimodal})}: A parameter-free fusion approach that serves as a stable baseline. It directly combines the visual and textual representations (e.g., via vector concatenation or element-wise averaging) without requiring additional learnable parameters:
    \begin{equation}
        \mathbf{x}_i = \mathbf{x}_i^{\text{vis}} \oplus \mathbf{x}_i^{\text{txt}},
    \end{equation}
    where \(\oplus\) denotes the concatenation operator.
    \item \textit{Weighted Attention Fusion (\texttt{multimodal\_attention})}: An adaptive fusion mechanism that dynamically learns attention weights to balance the contributions of visual and textual modalities on a per-item basis. This approach utilizes trainable projection layers and an attention gating network to automatically prioritize the most descriptive modality, yielding the highest recommendation accuracy:
    \begin{equation}
        \mathbf{x}_i = \beta_{v,i} W_v \mathbf{x}_i^{\text{vis}} + \beta_{t,i} W_t \mathbf{x}_i^{\text{txt}},
    \end{equation}
    where \(W_v \in \mathbb{R}^{d_m \times d_v}\) and \(W_t \in \mathbb{R}^{d_m \times d_t}\) are trainable projection matrices mapping visual and textual features to a shared \(d_m\)-dimensional space, and \(\beta_{v,i}\) and \(\beta_{t,i}\) are visual and textual attention coefficients satisfying \(\beta_{v,i} + \beta_{t,i} = 1\).
\end{itemize}

Advanced multimodal recommendation models integrate these fused features directly into their learning objectives:
\begin{itemize}
    \item \textit{BM3 \parencite{10.1145/3543507.3583251}}: Uses modality-specific projection layers and optimizes a self-supervised bootstrap contrastive loss to align user and item multimodal representations without requiring negative samples.
    \item \textit{FREEDOM}: Constructs a frozen item-item kNN similarity graph using item multimodal features, and propagates content embeddings over this static structure, alleviating the computational burden of dynamic graph convolutions on user-item bipartite graphs.
\end{itemize}

Despite the variety of multimodal fusion and graph-based models, there remains a lack of systematic evaluation comparing different fusion configurations (\texttt{img\_only}, \texttt{tfidf}, \texttt{multimodal}, \texttt{multimodal\_attention}) across different graph neural collaborative filtering architectures (LightGCN, CombiGCN, BM3, FREEDOM) in the fashion domain. This study addresses this gap by conducting a comprehensive comparative analysis of these modalities and fusion mechanisms under diverse structural propagation environments.
